\chapter{Algoritmo de cálculo del campo focal}
\label{apx:numerical_algorithm}

Los resultados numéricos presentados en el capítulo~\ref{ch:focal} se
obtuvieron con una implementación en Python de la formulación desarrollada en
esta tesis. Este apéndice describe el pseudocódigo esencial. Toda la
implementación reproducible se encuentra en el paquete \texttt{vecdiff}
distribuido junto a este trabajo; aquí se recogen únicamente los bloques
mínimos necesarios para reproducir las figuras.

\section{Discretización y estructuras}

Se emplean dos discretizaciones según la simetría del problema. Cuando el
campo incidente es axialmente simétrico se muestrea sobre una grilla radial
uniforme $\{\rho_i'\}_{i=0}^{N_\rho-1}$ en la pupila y $\{q_j\}_{j=0}^{N_q-1}$
en el plano focal, y se evalúa mediante transformadas de Hankel de orden
$m\in\{0,1,2\}$. Cuando el campo no posee simetría axial se muestrea sobre
una grilla cartesiana $\{(x_i,y_j)\}$ tanto en la pupila como en el plano
focal, y la propagación se realiza mediante FFT bidimensional. Ambas grillas
se convierten a la representación necesaria (cartesiana, cilíndrica o
circular) mediante los operadores de cambio de base descritos en el
capítulo~\ref{ch:focal}.

\section{Refracción en un dioptrio estigmático}

El núcleo del cálculo es la aplicación del operador de transmisión
$\mathbb T$ sobre la esfera de referencia objeto y su posterior propagación
al plano focal. El algoritmo~\ref{alg:single_diopter} lo resume para un
dioptrio.

\begin{algorithm}[H]
\caption{Campo focal tras un dioptrio estigmático}
\label{alg:single_diopter}
\begin{algorithmic}[1]
\Require Parámetros $(n_0,n_i,z_0,z_i,\lambda)$; superficie cartesiana
$z(\rho')$; campo incidente $\mathbb E_0(\rho',\varphi')$; modo
$\in\{\text{escalar},\text{vectorial}\}$.
\Ensure Campo focal $\mathbb E'(r,\varphi,f)$ y su intensidad total
$I(r,\varphi,f)$.
\State Muestrear la superficie: $z(\rho'_i),\;dz/d\rho'|_{\rho'_i}$
para cada nodo de pupila.
\State Calcular los cosenos de incidencia y transmisión desde la geometría
local y la ley de Snell.
\State Evaluar los coeficientes de Fresnel $t_p(\rho'_i)$, $t_s(\rho'_i)$.
\If{modo $=$ escalar}
  \State $t_p\gets 1,\;t_s\gets 1$ \Comment{ver
  \eqref{eq:t_plus_t_minus_def} y \eqref{eq:E_field_cartesian_paraxial}}
\EndIf
\State Construir el operador de transmisión en la base cartesiana
transversal:
$\mathbb T_{xy}=\tfrac{t_+}{2}\mathbb I_\perp+\tfrac{t_-}{2}\mathbb U(\phi')$,
según \eqref{eq:M_cartesian_tplus_tminus}.
\State Aplicar $\mathbb E^{\mathscr G'}\gets\mathbb T\,\mathbb E_0/\ell_0$ nodo a
nodo sobre la pupila.
\If{campo axialmente simétrico}
  \State Calcular $\mathbb E'_\perp$ mediante transformadas de Hankel de
  órdenes $\{0,1,2\}$, según la representación
  \eqref{eq:E_field_cartesian_paraxial}, \eqref{eq:E_L_E_R_paraxial} o
  \eqref{eq:E_cylindrical_paraxial}.
\Else
  \State Calcular $\mathbb E'_\perp$ mediante FFT bidimensional del campo
  sobre la pupila, con normalización de Fresnel.
\EndIf
\State Reconstruir la componente longitudinal $E'_z$ mediante la condición
de transversalidad
$K_z\,E'_z=-(K_x E'_x+K_y E'_y)$ en el espacio de frecuencias.
\State \Return $\mathbb E'=(E'_x,E'_y,E'_z)$ y $I=|\mathbb E'|^2$.
\end{algorithmic}
\end{algorithm}

El paso~5 aísla la aberración vectorial: fijar $t_p=t_s=1$ suprime el
canal $t_-$ y hace que $\mathbb T$ se reduzca a la identidad, con lo que el
campo focal colapsa sobre la predicción escalar de la difracción. Ambos
modos comparten la misma grilla y el mismo propagador, y sus resultados se
comparan directamente punto a punto.

\section{Composición de dos superficies}

El sistema litográfico del capítulo~\ref{ch:focal} está formado por dos
superficies cartesianas conjugadas separadas por un medio de sílice
fundida. El campo se propaga por composición de operadores: cada superficie
se trata mediante el algoritmo~\ref{alg:single_diopter}, y su salida se
utiliza como campo incidente sobre la superficie siguiente.

\begin{algorithm}[H]
\caption{Sistema de dos superficies cartesianas conjugadas}
\label{alg:two_diopter}
\begin{algorithmic}[1]
\Require Máscara de entrada $M(x,y)$; parámetros de cada superficie
$(n_0^{(k)},n_i^{(k)},z_0^{(k)},z_i^{(k)})$ para $k=1,2$; modo.
\Ensure Intensidad de imagen $I(x,y)$.
\State $\mathbb E_0\gets$ iluminación coherente de la máscara $M(x,y)$.
\State $\mathbb E_1\gets$ Algoritmo~\ref{alg:single_diopter} con parámetros
de la superficie~$1$ y campo incidente $\mathbb E_0$.
\State $\mathbb E_2\gets$ Algoritmo~\ref{alg:single_diopter} con parámetros
de la superficie~$2$ y campo incidente $\mathbb E_1$ (propagado hasta la
esfera de referencia de la superficie~$2$).
\State \Return $I=|\mathbb E_2|^2$.
\end{algorithmic}
\end{algorithm}

Este esquema se aplica dos veces —una en modo escalar, otra en modo
vectorial— sobre la misma máscara y con los mismos parámetros geométricos.
Los contrastes $C_H,\,C_V$ reportados en el capítulo~\ref{ch:focal} se
extraen de cortes horizontales y verticales sobre la imagen $I(x,y)$
resultante. La diferencia entre ambos modos aísla las manifestaciones de la
aberración vectorial estudiadas en esta tesis.

\section{Convergencia y verificaciones}

Los resultados se validaron mediante tres comprobaciones. Primero, la
consistencia entre las tres representaciones del capítulo~\ref{ch:focal}:
los mismos parámetros y las mismas condiciones de iluminación producen la
misma intensidad focal cuando se calculan en base cartesiana, cilíndrica y
circular. Segundo, la reducción al límite escalar: fijando
$t_p=t_s=1$ y anulando la reconstrucción longitudinal, el algoritmo
reproduce el patrón de Airy hasta la precisión de la grilla, verificando
que la modulación por $\mathbb T$ es la única fuente de la desviación
vectorial. Tercero, el refinamiento de malla: los contrastes
$C_H,\,C_V$ del ejemplo litográfico, así como los radios principales
$R_\parallel,\,R_\perp$ del punto focal deformado, se calcularon sobre
grillas cada vez más finas hasta que sus variaciones relativas quedaron por
debajo del $0.5\,\%$; los valores reportados en el capítulo~\ref{ch:focal}
corresponden a la primera grilla que satisface esta tolerancia. La misma
tolerancia se usó para fijar el orden de truncamiento de las series
paraxiales del apéndice~\ref{apx:ovoids_series}: los coeficientes de
Fresnel se retuvieron hasta el orden a partir del cual la inclusión del
siguiente cambiaba los observables medidos en menos de esa fracción.
